% !TEX program = lualatex
\documentclass[aspectratio=43,professionalfonts,handout]{beamer}
\usepackage{beamer_template}

\newcommand{\LectureNum}{8}
\newcommand{\LectureTitle}{周期2πの周期関数のフーリエ級数}
\newcommand{\TextPages}{p.\,74-79}
\newcommand{\CourseName}{応用数学}
\renewcommand{\TermName}{前期}

\begin{document}

\begin{frame}{\CourseName\quad 第\LectureNum 回「\LectureTitle」}
  {\small \TermName\quad \TeacherName\quad ／\quad 教科書 \TextPages}

  \textbf{目標}：周期2πの周期関数のフーリエ係数を求め、フーリエ級数を計算できる

  \begin{block}{概要}
  \begin{itemize}
    \item フーリエ級数とは
    \item フーリエ係数
    \item 収束定理
  \end{itemize}
  \end{block}
\end{frame}

\begin{frame}{フーリエ級数の概念}
  \begin{block}{}
  \begin{itemize}
    \item 任意の周期関数を三角関数の和で表す
    \item 音の波形の例：ピアノvsフルート
    \item 工学・物理での重要性
  \end{itemize}
  \end{block}
\end{frame}

\begin{frame}{直交性と係数の公式}
  \begin{block}{}
  \begin{itemize}
    \item 三角関数の直交性（I〜IV）
    \item $c_{0}, a_{n}, b_{n}$ の導出
    \item フーリエ係数の計算手順
  \end{itemize}
  \end{block}
\end{frame}

\begin{frame}{例題1の計算}
  \begin{block}{}
  \begin{itemize}
    \item $f(x)$ の設定とグラフ
    \item $c_{0}, a_{n}, b_{n}$ を順に計算
    \item フーリエ級数の表示
  \end{itemize}
  \end{block}
\end{frame}

\begin{frame}{収束定理とギブス現象}
  \begin{block}{}
  \begin{itemize}
    \item 区分的に滑らかな関数への収束
    \item 不連続点での値 $(f(x+0)+f(x-0))/2$
    \item ギブス現象：$N \to \infty$ でも消えない振動
  \end{itemize}
  \end{block}
\end{frame}

\begin{frame}{演習}
  \begin{exampleblock}{自分で解いてみよう}
  \begin{itemize}
    \item 問2〜5より選題
  \end{itemize}
  \end{exampleblock}
\end{frame}

\begin{frame}{まとめ}
  \begin{itemize}
    \item フーリエ係数の公式
    \item 収束定理
    \item 次回：一般の周期関数
  \end{itemize}
\end{frame}

\end{document}
